% ═════════════════════════════════════════════════════════════════════════════
%  第二章  相关工作
% ═════════════════════════════════════════════════════════════════════════════
\section{相关工作}
\label{sec:related}

% ─────────────────────────────────────────────────────────────────────────────
\subsection{检索增强生成（RAG）综述}
\label{sec:related-rag}

\subsubsection{RAG技术演进}

检索增强生成（Retrieval-Augmented Generation，RAG）由Lewis等人\cite{lewis2020rag}
提出，其核心思想是在LLM推理阶段动态检索外部知识库，
将召回的相关文本片段注入上下文窗口，引导模型生成更准确、更具事实依据的答案。
Fan等人\cite{fan2024ragsurvey}从架构设计、训练策略与应用场景三个维度
对RAG-LLMs研究进行了全面梳理，认为RAG在开放域问答等知识密集型任务中展现出显著潜力。
Gao等人\cite{gao2023ragsurvey}的综述将RAG的发展划分为三个递进阶段：

\begin{figure}[htbp]
  \centering
  \includegraphics[width=\linewidth]{figures/fig7_rag_evolution.pdf}
  \caption{RAG技术三阶段演进对比：Naive RAG、Advanced RAG与GraphRAG}
  \label{fig:rag-evolution}
\end{figure}

\textbf{Naive RAG}是最早期范式，以稠密向量检索为核心，
直接将Top-$K$文档块拼接入提示词。
该范式在通用问答任务上表现有效，
但在专业领域面临检索精度低、上下文噪声大等问题\cite{fan2024ragsurvey}。
Liu等人\cite{liu2023lostinmiddle}通过实验发现，
当无关段落混入检索上下文时，LLM的生成质量会显著下降——
这一"中间丢失"（Lost in the Middle）现象在长上下文场景尤为突出，
揭示了Naive RAG在精确查询场景中的根本局限。

\textbf{Advanced RAG}在上述基础上引入查询改写、混合检索与精排等优化模块，
以多层次干预系统性提升检索质量。
代表性改进包括：
多查询扩展\cite{ma2023query}（用LLM将原始查询改写为多种表述后并行检索）、
稀疏-稠密混合检索\cite{luan2021sparse}（BM25与向量检索互补融合）、
以及交叉编码器精排\cite{nogueira2019passage}（在粗排后进行细粒度语义重排）。

\textbf{GraphRAG}则走向了更深层的结构化方向。
Edge等人\cite{edge2024graphrag}提出将知识图谱与RAG结合，
以图结构表示实体关系，支持多跳推理，
有效弥补了向量检索在结构化推理方面的能力边界。
本文在GraphRAG框架的基础上，针对航空装配场景的特殊需求，
进一步设计了三源知识图谱与基于密集语义近邻搜索（ANN）的子图召回路径——
通过将BGE-M3统一编码用户查询与图谱实体描述，使中英混合查询能够直接命中
图谱中的语义近邻锚点，再以一跳邻域扩展提取相关子图，
解决了纯文本图谱无法覆盖几何约束知识、且关键词匹配对中文查询召回率偏低的双重问题。

\subsubsection{稀疏-稠密混合检索}

早期RAG工作以双编码器（Bi-Encoder）稠密检索为主，
代表性模型包括DPR\cite{karpukhin2020dpr}和BGE系列\cite{bge2023}。
稠密检索擅长捕捉语义相似性，但在面对零件编号（PN）等精确标识符时匹配能力不足。

BM25\cite{robertson2009bm25}作为经典概率检索模型，
在词汇层面具有天然的精确匹配优势。
研究表明，稀疏检索与稠密向量检索互补融合后，
综合召回率和准确率通常优于任一单路方法\cite{luan2021sparse}。
现有工业RAG工作（如HybridRAG\cite{shen2025hybridrag}）采用
线性加权融合（$\text{score} = \alpha \cdot s_{\text{dense}} + \beta \cdot s_{\text{sparse}}$），
该方式依赖人工调整权重且对评分分布敏感。
本文改用互惠排名融合（RRF）\cite{cormack2009rrf}，
无需对不同评分空间进行归一化处理，对单路检索失败具有天然鲁棒性，
且可自然地将文本块与图谱子图等异质内容统一排序。

精排阶段，交叉编码器（Cross-Encoder）基于Transformer编码器
\cite{vaswani2017attention,devlin2019bert}，
在推理时对查询与候选段落进行全注意力交互，
相比双编码器的独立编码范式能更精准地捕捉细粒度语义对齐关系\cite{nogueira2019passage}。

% ─────────────────────────────────────────────────────────────────────────────
\subsection{知识图谱构建与推理}
\label{sec:related-kg}

知识图谱以三元组~$(e_h, r, e_t)$~的形式表示实体间结构化关系，
为需要多跳推理的复杂问答提供向量检索难以胜任的能力。
Hogan等人\cite{hogan2021kg}对知识图谱的表示方法、构建技术与应用进行了权威综述，
指出KG在知识密集型任务中具有重要的结构化推理优势。
Ji等人\cite{ji2022kgsurvey}从表示学习、知识获取与应用三个维度系统梳理了KG研究进展，
为本文的本体设计和实体对齐方法奠定了理论基础。
Pan等人\cite{pan2024kgllmsurvey}进一步梳理了知识图谱与LLM融合的路线图，
认为两者结合在专业领域问答中具有显著协同优势——
这也是本文选择KG+RAG融合路线的核心依据之一。

\subsubsection{基于LLM的零样本三元组抽取}

早期命名实体识别（NER）与关系抽取（RE）工作依赖大规模标注数据，
在航空制造这类低资源领域面临标注成本高昂的瓶颈。
近年来，Pan 等\cite{pan2024kgllmsurvey}的综述系统梳理了
基于预训练语言模型的零样本或少样本三元组抽取范式——
通过将本体模式编码进少样本提示词（Few-shot Prompting）驱动 LLM 直接输出结构化三元组——
可大幅降低知识图谱的构建成本，使低资源领域的 KG 构建成为可能。

为进一步提升抽取质量，Edge等人\cite{edge2024graphrag}提出的Gleaning机制
在首轮抽取后以第二轮提示要求模型补充遗漏实体；
实体对齐（Entity Linking）则通过缩写词典映射、编辑距离模糊匹配和语义向量
三级级联方式\cite{wu2020corefresolution}解决同义表达冲突，
从而保证图谱内部的一致性。

\subsubsection{基于CAD几何数据的结构化知识抽取}

现有工业知识图谱研究多聚焦于文本、数据库和传感器数据，
鲜有工作直接从三维CAD模型中抽取精确的结构化知识。
Shi等人\cite{shi2022arkg}提出基于"功能-结构-工序-资源"多维本体的
装配资源知识图谱（ARKG），通过本体映射将CAD模型的实例化数据融入知识图谱，
实现了产品、工序和装配资源信息的有效耦合；
但该工作主要面向CAD属性数据，未涉及几何约束的系统性解析。

事实上，STEP/ISO 10303\cite{srinivasan2017bom}格式的CAD文件以标准化方式编码了
零件的父子包含关系、几何配合约束（同轴、贴合、平行等）和制造特征（PMI）——
这些信息以确定性方式存在，精度高且无需人工标注。
本文提出的CAD解析轨道正是利用OCCT从STEP文件中
提取装配层级树（\textit{isPartOf}）、配合约束（\textit{matesWith}）
与空间邻接关系（\textit{adjacentTo}），
弥补了纯文本知识图谱在几何拓扑知识方面的空白，
构成本文相较于现有工业 RAG 工作的主要差异之一。

\subsubsection{BOM驱动的实体注册与对齐}

BOM（Bill of Materials）是制造业中零件身份的权威来源，
包含零件编号（PN）、名称、数量和材料规格等结构化属性\cite{srinivasan2017bom}。
Liu等人\cite{liu2023avkg}指出，
航空装配领域的数据源高度非结构化，实体边界模糊
（如"雷达校准"可归类为操作实体，也可将"雷达"单独归为部件实体），
这一问题在跨源实体对齐时尤为突出。
以BOM作为图谱实体节点的注册基础，可有效解决多源数据中零件名称不一致的问题：
手册中的"主轴承"经实体链接与BOM中的 PN:~1002-AX 对应，
CAD中的几何体引用亦通过STEP文件的\texttt{PRODUCT}实体表解析后与BOM对齐，
从而在三源异构数据之间建立统一的实体标识空间。

\subsubsection{多智能体知识抽取}

由于三类数据源的格式与语义差异显著，单一流水线难以高质量处理全部来源。
Guo等人\cite{guo2024multiagent}证明，多智能体系统（Multi-Agent System）
通过分工协同——各智能体专注于自身最擅长的数据类型——总体性能优于单体方案。
Park等人\cite{park2023generative}进一步展示了多智能体协作在复杂推理任务中的潜力。
受此启发，本文设计了BOM解析智能体、文本抽取智能体、
CAD结构化解析智能体与验证智能体等专职智能体，
协同完成三源知识的抽取、融合与一致性校验。

\subsubsection{图谱推理与RAG融合}

GraphRAG\cite{edge2024graphrag}通过识别命名实体为锚节点、可变深度邻域扩展完成多跳遍历，
有效支持"某零件的所有子装配体"等结构化查询，
但其图谱仅来源于文本抽取，不含几何信息；
此外，传统 GraphRAG 多依赖关键词字符串匹配定位锚节点，
对中英混合查询、近义术语场景的鲁棒性有限。
本文的知识图谱检索路径（路径三）在此基础上做出两项改进：
其一，以CAD解析所得的几何约束节点丰富了图谱的结构深度，
从而支持纯文本GraphRAG无法回答的几何拓扑类查询；
其二，将锚节点定位从关键词CONTAINS匹配改为基于BGE-M3向量的密集语义近邻搜索，
使中英混合查询在统一编码空间内直接命中相关实体，再以一跳邻域扩展形成上下文子图。
这一设计减少了对精确字符串命中的依赖，为跨章节、跨语言别名查询提供了更稳定的锚节点定位方式。

% ─────────────────────────────────────────────────────────────────────────────
\subsection{查询扩展与生成增强}
\label{sec:related-generation}

\subsubsection{多查询扩展}

多查询扩展（Multi-Query Expansion）是缓解词汇不匹配的常用技术。
Ma等人\cite{ma2023query}证明，
利用LLM将原始查询改写为多种不同表述后并行检索、
再对结果合并去重，可显著提升召回覆盖率。
本文将多查询扩展应用于全部三条检索路径，
构成查询集~$\mathcal{Q} = \{q, q_1, q_2\}$，
是本系统实现高召回率的重要技术手段之一。

% ─────────────────────────────────────────────────────────────────────────────
\subsection{航空发动机装配智能化}
\label{sec:related-aero}

航空发动机装配是制造业中精度要求极高、工艺高度复杂的环节之一。
传统装配辅助系统多依赖人工编写的结构化数据库和固定检索规则，
缺乏对自然语言查询的理解能力，在面对机型更新或工艺变更时适应性较差。
近年来，随着知识图谱与大语言模型技术的快速发展，
该领域的智能化研究日趋活跃，可从以下三个维度加以梳理。

\textbf{航空领域知识图谱应用。}
Tang等人\cite{tang2023aircraft}系统综述了航空故障诊断知识图谱的构建与应用研究；
Meng等人\cite{meng2024aircraft}构建了面向航空PHM（预测与健康管理）的故障知识图谱平台，
涵盖航空发动机典型故障的实体建模和推理路径；
Wu等人\cite{wu2024engine}专门构建了航空发动机润滑系统故障知识图谱，
并在实际工程场景中验证了KG辅助故障诊断的有效性。
上述工作一致表明，专业化知识图谱是提升航空领域故障诊断精度的关键手段；
然而它们均聚焦于故障诊断场景，
未涉及装配引导这一需要三源异构知识深度融合的精细任务。

\textbf{航空装配知识图谱构建。}
Liu等人\cite{liu2023avkg}利用联合知识抽取模型自动从装配手册文本中
提取实体和关系，构建了含5类实体（步骤、操作、工具、部件、设施）和
3类关系（Component-Whole、Content-In、Instrument-Agency）的航空装配KG，
以Neo4j作为图数据库存储，实体和关系抽取的F1-score分别达89.71\%和91.27\%；
Shi等人\cite{shi2022arkg}构建了基于"功能-结构-工序-装配资源"多维本体的
装配资源知识图谱，将CAD实例数据和工程知识在统一语义框架下集成，
为装配工艺设计提供知识共享与复用支持。
两项工作均验证了KG在装配领域的可行性，
但其数据融合粒度仍有限，
尚未将BOM身份注册、技术手册工序知识与三维几何约束纳入统一图谱框架。

\textbf{LLM与装配智能化融合。}
Blumenthal等人\cite{blumenthal2020nlp}探索了将技术手册文本结构化以支持关键字检索的方案；
Zhao等人\cite{zhao2022assembly}提出基于BERT\cite{devlin2019bert}的装配指令理解模型，
在有监督条件下提升了指令语义解析能力。
然而，上述工作均未引入大语言模型的生成能力，难以应对表述多样的开放域问题。
Zhao等人\cite{zhao2023llmsurvey}的综述表明，LLM在自然语言理解、
推理和生成方面具有显著的涌现能力；
Wu 等\cite{wu2024engine}将 LLM 与航空发动机润滑系统 KG 协同，
证明了 LLM 在装备级专业故障问答中的应用潜力，
但纯LLM方案存在幻觉和知识陈旧问题，
在涉及精确数值（扭矩、公差、零件编号）的查询上可靠性不足\cite{fan2024ragsurvey}。

Liu等人\cite{liu2024jointkg}则进一步将航空装配知识图谱通过前缀微调嵌入LLM，
利用2-step子图驱动前缀调优，在功能测试场景下实现了98.5\%的故障诊断准确率，
验证了KG结构化知识与LLM推理能力深度融合的可行性。

\textbf{与本文最相关的工作。}
Shen等人\cite{shen2025hybridrag}的HybridRAG
在飞机故障排查场景下证明了KG+向量混合检索的有效性（F1提升4\%、幻觉率降低7\%），
并设计了Planner-Agent迭代探索算法，通过计划-执行-自我反思的循环提升答案质量。
本文借鉴其KG增强检索思想，并针对航空发动机装配的特殊需求做出三项关键扩展：
（1）引入三源（BOM+CAD+手册）多智能体知识图谱构建流水线，
将本体规模从手册文本的5类实体扩展至7类实体和9类关系，
构建覆盖更完整的装配知识底座；
（2）将KG锚节点定位从关键词匹配改为BGE-M3语义近邻搜索，
提升中英混合查询与专业别名查询下的图谱召回稳定性；
（3）将线性加权融合改进为RRF无参融合，提升文本块与图谱子图融合的鲁棒性。
上述扩展使本系统在知识覆盖度、结构化推理与融合可靠性上
相对现有工作具有差异化优势，并直接面向装配引导这一高精度应用场景。
